
\documentclass[11pt,a4paper]{article}
\usepackage[T1]{fontenc}
\usepackage{lmodern}
\usepackage{amsmath,amssymb,amsfonts,mathtools}
\usepackage{graphicx,booktabs,longtable,array,tabularx,adjustbox}
\usepackage{geometry,microtype,fancyhdr,lastpage,etoolbox,enumitem}
\usepackage{newunicodechar,textcomp}
\usepackage{hyperref}
\hypersetup{hidelinks,pdfcreator={XeLaTeX / QM1 academic edition}}
\geometry{left=20mm,right=20mm,top=22mm,bottom=23mm}
\setcounter{secnumdepth}{2}
\setlength{\parindent}{0pt}
\setlength{\parskip}{4pt plus 1pt}
\setlength{\emergencystretch}{2em}
\setlength{\tabcolsep}{3pt}
\renewcommand{\arraystretch}{1.16}
\setlist{nosep,topsep=3pt,itemsep=2pt,leftmargin=7mm}
\AtBeginEnvironment{longtable}{\footnotesize\setlength{\tabcolsep}{2.2pt}}
\newcommand{\tightlist}{\setlength{\itemsep}{0pt}\setlength{\parskip}{0pt}}
\providecommand{\Cov}{\operatorname{Cov}}
\providecommand{\Var}{\operatorname{Var}}
\providecommand{\Corr}{\operatorname{Corr}}
\providecommand{\MAD}{\operatorname{MAD}}
\providecommand{\E}{\mathbb{E}}
\providecommand{\R}{\mathbb{R}}
\newunicodechar{—}{\textemdash}
\newunicodechar{–}{\textendash}
\newunicodechar{€}{\texteuro}
\newunicodechar{≤}{\ensuremath{\leq}}
\newunicodechar{≥}{\ensuremath{\geq}}
\newunicodechar{×}{\ensuremath{\times}}
\newunicodechar{÷}{\ensuremath{\div}}
\newunicodechar{Σ}{\ensuremath{\Sigma}}
\newunicodechar{σ}{\ensuremath{\sigma}}
\newunicodechar{μ}{\ensuremath{\mu}}
\newunicodechar{Δ}{\ensuremath{\Delta}}
\newunicodechar{∑}{\ensuremath{\sum}}
\newunicodechar{≈}{\ensuremath{\approx}}
\newunicodechar{√}{\ensuremath{\surd}}
\newunicodechar{≠}{\ensuremath{\ne}}
\newunicodechar{→}{\ensuremath{\to}}
\newunicodechar{∞}{\ensuremath{\infty}}
\newunicodechar{∈}{\ensuremath{\in}}
\newunicodechar{−}{\ensuremath{-}}
\newunicodechar{²}{\ensuremath{^{2}}}
\newunicodechar{³}{\ensuremath{^{3}}}
\newunicodechar{₀}{\ensuremath{_{0}}}
\newunicodechar{₁}{\ensuremath{_{1}}}
\newunicodechar{₂}{\ensuremath{_{2}}}
\newunicodechar{₃}{\ensuremath{_{3}}}
\newunicodechar{₄}{\ensuremath{_{4}}}
\newunicodechar{ₜ}{\ensuremath{_{t}}}
\newunicodechar{ₚ}{\ensuremath{_{p}}}
\newunicodechar{ₙ}{\ensuremath{_{n}}}
\newunicodechar{ₓ}{\ensuremath{_{x}}}
\newunicodechar{ᵢ}{\ensuremath{_{i}}}
\newunicodechar{ŷ}{\ensuremath{\hat y}}
\newunicodechar{Ŷ}{\ensuremath{\widehat Y}}
\newunicodechar{ₘ}{\ensuremath{_{m}}}
\newunicodechar{ₑ}{\ensuremath{_{e}}}
\newunicodechar{ₐ}{\ensuremath{_{a}}}
\newunicodechar{ₖ}{\ensuremath{_{k}}}
\newunicodechar{ₛ}{\ensuremath{_{s}}}
\newunicodechar{ₗ}{\ensuremath{_{l}}}
\newunicodechar{ₕ}{\ensuremath{_{h}}}
\newunicodechar{ₒ}{\ensuremath{_{o}}}
\newunicodechar{₊}{\ensuremath{_{+}}}
\newunicodechar{₋}{\ensuremath{_{-}}}
\newunicodechar{₌}{\ensuremath{_{=}}}
\newunicodechar{₍}{\ensuremath{_{(}}}
\newunicodechar{₎}{\ensuremath{_{)}}}
\newunicodechar{⁻}{\ensuremath{^{-}}}
\newunicodechar{⁺}{\ensuremath{^{+}}}
\newunicodechar{⁰}{\ensuremath{^{0}}}
\newunicodechar{¹}{\ensuremath{^{1}}}
\newunicodechar{⁴}{\ensuremath{^{4}}}
\newunicodechar{⁵}{\ensuremath{^{5}}}
\newunicodechar{⁶}{\ensuremath{^{6}}}
\newunicodechar{⁷}{\ensuremath{^{7}}}
\newunicodechar{⁸}{\ensuremath{^{8}}}
\newunicodechar{⁹}{\ensuremath{^{9}}}
\makeatletter
\renewcommand{\section}{\@startsection{section}{1}{0pt}{2.0ex plus .7ex minus .3ex}{1.0ex}{\large\bfseries}}
\renewcommand{\subsection}{\@startsection{subsection}{2}{0pt}{1.8ex plus .5ex minus .2ex}{.7ex}{\normalsize\bfseries}}
\renewcommand{\subsubsection}{\@startsection{subsubsection}{3}{0pt}{1.3ex plus .4ex minus .2ex}{.5ex}{\normalsize\bfseries}}
\makeatother
\pagestyle{fancy}
\fancyhf{}
\renewcommand{\headrulewidth}{0.35pt}
\renewcommand{\footrulewidth}{0pt}
\fancyhead[L]{\small\textsc{KEDGE Business School --- QM1}}
\fancyhead[R]{\small\textsc{2026--2027}}
\fancyfoot[L]{\small Badr ABOUFARASSE, PhD}
\fancyfoot[R]{\small\thepage\ /\ \pageref{LastPage}}
\setlength{\headheight}{14pt}
\setkeys{Gin}{keepaspectratio}
\begin{document}

\begin{center}
{\small\textsc{KEDGE Business School \quad EBP1 \quad Quantitative Methods 1}}\\[6pt]
{\Large\bfseries Worksheet 5 — Time series}\\[6pt]
{\normalsize\bfseries Exercise statements}\\[6pt]
{\small Badr ABOUFARASSE, PhD \quad --- \quad 2026--2027}
\end{center}
\vspace{2mm}\hrule\vspace{6mm}

\hypertarget{core-applications}{%
\section{Core applications}\label{core-applications}}

Master the essential definitions and calculations from the course.

\emph{Core application --- 10 min}

\hypertarget{exercise-1-components-and-models}{%
\subsection{Exercise 1 -- Components and models}\label{exercise-1-components-and-models}}

\begin{enumerate}
\item
  Match long-run trend, recession, recurring Q4 peak and exceptional shock to {\(T_t\)}, {\(C_t\)}, {\(S_t\)} or {\(R_t\)}.
\item
  Write the full additive and multiplicative models.
\item
  Explain when multiplicative seasonality is more plausible than additive seasonality.
\end{enumerate}

\emph{Core application --- 20 min}

\hypertarget{exercise-2-moving-averages-and-centring}{%
\subsection{Exercise 2 -- Moving averages and centring}\label{exercise-2-moving-averages-and-centring}}

Quarterly demand over two years is {\[80,\
100,\ 120,\ 140,\ 88,\ 110,\ 132,\ 154.\]} Describe the main pattern, calculate the five 4-quarter moving averages, centre them for periods 3 to 6, and explain why an even-order moving average must be centred.

\hypertarget{assessment-and-exam-practice}{%
\section{Assessment and exam practice}\label{assessment-and-exam-practice}}

Practise independently with assessment and final-exam style questions.

\emph{Assessment / exam practice --- 25 min}

\hypertarget{exercise-3-multiplicative-decomposition-and-forecasting}{%
\subsection{Exercise 3 -- Multiplicative decomposition and forecasting}\label{exercise-3-multiplicative-decomposition-and-forecasting}}

\begin{longtable}[]{@{}lllll@{}}
\toprule\noalign{}
Year & Q1 & Q2 & Q3 & Q4 \\
\midrule\noalign{}
\endhead
\bottomrule\noalign{}
\endlastfoot
2005 & 82 & 98 & 120 & 150 \\
2006 & 91 & 111 & 137 & 171 \\
2007 & 104 & 126 & 154 & 194 \\
\end{longtable}

Given {\(\widehat T_t=6.343t+86.939\)} and seasonal coefficients {\(0.789,0.903,1.052,1.255\)}, interpret the trend slope, check normalization, forecast the four quarters of 2008, and explain multiplicative deseasonalization.

\emph{Assessment / exam practice --- 65 min}

\hypertarget{exercise-5-full-seasonal-decomposition-exam-style}{%
\subsection{Exercise 5 --- Full seasonal decomposition --- exam style}\label{exercise-5-full-seasonal-decomposition-exam-style}}

A firm observes quarterly sales over three years. Use the multiplicative model Yₜ=Tₜ×Sₜ, without explicit cycle or irregular component. Do not round intermediate values.

\begin{longtable}[]{@{}lllll@{}}
\toprule\noalign{}
Année / Year & Q1 & Q2 & Q3 & Q4 \\
\midrule\noalign{}
\endhead
\bottomrule\noalign{}
\endlastfoot
2019 & 40 & 52 & 48 & 68 \\
2020 & 44 & 58 & 53 & 76 \\
2021 & 50 & 65 & 60 & 84 \\
\end{longtable}

\begin{enumerate}
\tightlist
\item
  Plot the series and calculate all nine 4-quarter moving averages.
\item
  Centre the moving averages (t=3 to 10). Build the t, CMA4, t², t×CMA4 table.
\item
  Estimate trend Tₜ=at+b by OLS on the eight centred moving averages.
\item
  Calculate Yₜ/Tₜ for all 12 quarters, the four seasonal means and normalised factors summing to 4.
\item
  Forecast all four quarters of 2022 (t=13 to 16), first deseasonalised, then seasonalised.
\end{enumerate}

\hypertarget{optional-advanced-practice}{%
\section{Optional advanced practice}\label{optional-advanced-practice}}

Explore additional interpretations and professional applications.

\emph{Advanced practice --- 15 min}

\hypertarget{exercise-4-forecast-evaluation}{%
\subsection{Exercise 4 -- Forecast evaluation}\label{exercise-4-forecast-evaluation}}

For this exercise, suppose actual 2008 values were {\(136,162,188,245\)}. Using the Exercise 3 forecasts, calculate forecast errors, MAE and RMSE, and explain the value of a separate validation period and a naive benchmark.

\emph{Advanced practice --- 20 min}

\hypertarget{extra-1-forecast-validation-and-structural-break}{%
\subsection{Extra 1 --- Forecast validation and structural break}\label{extra-1-forecast-validation-and-structural-break}}

Actual monthly demand is 100, 104, 108, 150, 154 and 158. A trend model fitted on the first three months predicts 100, 104, 108, 112, 116 and 120.

\begin{enumerate}
\tightlist
\item
  Calculate forecast errors for months 4--6.
\item
  Calculate MAE and mean error over the validation period.
\item
  Diagnose the model.
\item
  Propose a responsible forecasting response.
\end{enumerate}

\textbf{Method hint:}

Use error = actual − forecast. A persistently positive error signals systematic underforecasting.

\end{document}
